% !TeX program = xelatex
% 编译：xelatex 典型的信号.tex（VS Code 的 LaTeX Workshop 会读上一行，自动用 xelatex）
% 来源：2022 浙江大学《信号与系统》胡浩基老师，第 1 章前两节课（视频约 59 分钟）。
\documentclass[UTF8,a4paper,11pt]{ctexart}
\usepackage{amsmath,amssymb}
\usepackage{geometry}
\geometry{margin=2.4cm}
\usepackage{xcolor}
\usepackage{mdframed}
\usepackage{booktabs}
\usepackage{enumitem}

\setCJKmainfont{Songti SC}[BoldFont={Songti SC Bold}]
\setCJKsansfont{PingFang SC}

\newmdenv[backgroundcolor=blue!4,linecolor=blue!50!black,linewidth=1pt,innertopmargin=6pt,innerbottommargin=6pt]{rememberbox}
\newenvironment{remember}{\begin{rememberbox}\textbf{必须背下来}\par\smallskip}{\end{rememberbox}}
\newmdenv[backgroundcolor=yellow!8,linecolor=orange!70!black,linewidth=1pt,innertopmargin=6pt,innerbottommargin=6pt]{notebox}
\newenvironment{note}{\begin{notebox}\textbf{说明}\par\smallskip}{\end{notebox}}
\newcommand{\Sa}{\operatorname{Sa}}
\renewcommand{\d}{\,\mathrm{d}}

\title{典型的信号\\\large 浙江大学 2022《信号与系统》（胡浩基）第 1 章 听课笔记}
\author{}
\date{}

\begin{document}
\maketitle

\begin{note}
这份笔记由视频整理而来，来源要分清：
\begin{itemize}[nosep]
\item 概念与口述内容：来自语音转写。转写模型很小，错字很多（如「阶跃」被写成「接越」），我按上下文还原；
\item 板书公式与证明：来自画面。第一遍转写的 41:30--54:30 一段因 whisper 复读而失效，后来换用 large-v3-turbo 模型（关闭上下文、加术语提示词）对 40:00--56:00 重新转写，这一段已基本恢复；证明的步骤我用重转写的讲解与板书截图互相对照过；
\item 标注「（补充）」的内容是我自己加的，课上没讲，或讲了但我没听清。
\end{itemize}
\end{note}

\section{一句话主干}
本节讲的是：\textbf{信号怎么分类}，以及\textbf{后面整门课都会反复用到的四个基本信号}——阶跃(unit step)、冲激(impulse)、抽样函数(sampling function)，外加离散里对应的单位脉冲和单位阶跃。其中有两个「必须记住的积分结果」，后面推傅里叶变换要用。

%==================================================================
\section{信号与系统是什么}
\begin{itemize}
\item \textbf{信号}(signal)：表达、传递信息的符号。
\item \textbf{系统}(system)：有输入、有输出。接收一个输入信号，处理之后变成另一个信号输出。
\end{itemize}

%==================================================================
\section{信号的分类}
老师说：书上第 3 页起有很多分类，只讲重要的，有时还会自己加一个。第一个「按维度分」就是他自己加的，不在书上。

\subsection{按维度分：这门课只讨论一维}
\begin{tabular}{@{}lll@{}}
\toprule
维度 & 例子 & 自变量 $\to$ 取值 \\
\midrule
一维 & 声音（语音）信号 & 时间 $\to$ 声音强度 \\
二维 & 图像 & （横坐标，纵坐标）$\to$ RGB 三个数 \\
三维 & 视频；深度图 & 图像两个维度 $+$ 时间；图像两个维度 $+$ 深度 \\
四维 & VR 眼镜看到的画面 & 三维 $+$ 时间 \\
\bottomrule
\end{tabular}

\medskip
老师的几个说法：
\begin{itemize}[nosep]
\item 图像里每个点用 RGB 三个数编码，按课上的说法每个数用 24 比特（取值 $0\sim 2^{24}-1$）。\emph{（补充：两次独立转写（base 与 large-v3-turbo）都是这个说法，所以这是老师原话，不是听错。但与常见事实不符：常见做法是红、绿、蓝\textbf{三个通道各 8 比特，合起来一个像素 24 比特}，取值 $0\sim255$。老师这里很可能是把「一个像素的总比特数」说成了「每个数的比特数」。这属于口误，对后面的内容没有影响，以 8 比特/通道为准，我没有核对书。）}
\item 人看到一幅图像，计算机看到的是三个矩阵（红、绿、蓝各一个）。
\item 视频的帧率：传统 24 帧，高清 50 帧。
\item \textbf{只讨论一维}，原因是简单：一维是学二、三、四维的基础，不是高维不重要，而是它们在别的课里学。二维傅里叶变换会给公式，但不展开。
\end{itemize}

\subsection{连续信号与离散信号}
\begin{remember}
\begin{center}
圆括号 $x(t)$ $\Rightarrow$ 连续信号，$t\in\mathbb{R}$\qquad\qquad
方括号 $x[n]$ $\Rightarrow$ 离散信号，$n\in\mathbb{Z}$
\end{center}
整本书统一这个写法，\textbf{没有 $x[t]$，也没有 $x(n)$}。老师特别提醒：学到后面做离散傅里叶变换，明明写的是方括号，做着做着套回连续的公式，是常见的错。
\end{remember}

学生问过：为什么 $n$ 只能取整数，不能取 $0.5$？老师的回答：$n$ 只是一个「名字」，只表示「我不在实数轴上取值，而在一个离散的集合上取值」；用 $-0.5,0.5$ 当名字也行，但为了统一，约定用整数。

\subsection{周期信号与非周期信号}
\textbf{连续}：
\[
x(t)=x(t+mT),\quad m\in\mathbb{Z}
\]
$T$ 取\textbf{最小正周期}。若 $T$ 是周期，则 $2T,3T,\dots$ 也都是周期，其中最小的正的那个叫基本周期。

\textbf{离散}：
\[
x[n]=x[n+mN],\quad m\in\mathbb{Z},\ N\ \text{为正整数}
\]
老师口述这里的 $N$ 是「自然数」，我写成正整数是因为 $N=0$ 时「重复」没有意义。老师的教学要求：「举一反三、触类旁通」。以后他有时只给连续或只给离散的定义，另一个要自己补出来。

\subsection{奇信号与偶信号}
\begin{align*}
\text{偶信号}\ (\text{even}):&\quad x(t)=x(-t)\quad \text{关于纵轴对称}\\
\text{奇信号}\ (\text{odd}):&\quad x(t)=-x(-t)\quad \text{关于原点旋转 }180^\circ\text{ 对称}
\end{align*}
老师评价：这就是高中的奇函数、偶函数，不算本课新知识。

\subsubsection*{本课第一个真正的新知识：任何信号都能唯一分解成「偶 $+$ 奇」}
\begin{remember}
\textbf{定理。}对任意信号 $x(t)$（定义域关于原点对称，如 $t\in\mathbb{R}$），存在\textbf{唯一}的一对信号 $x_e(t)$、$x_o(t)$，其中 $x_e$ 是偶信号、$x_o$ 是奇信号，使得
\[
x(t)=x_e(t)+x_o(t),\qquad\text{并且}\quad
x_e(t)=\frac{x(t)+x(-t)}{2},\quad x_o(t)=\frac{x(t)-x(-t)}{2}.
\]
\end{remember}
这里 $e$ 是 even，$o$ 是 odd，下标是这么来的。要证明的有两件事：\textbf{存在}（至少有一种拆法）和\textbf{唯一}（至多有一种拆法）。

\paragraph{证明一：存在性。}直接把上面的 $x_e,x_o$ 定义出来，验三条：
\begin{enumerate}[nosep]
\item \textbf{$x_e$ 是偶的}：
$x_e(-t)=\dfrac{x(-t)+x\bigl(-(-t)\bigr)}{2}=\dfrac{x(-t)+x(t)}{2}=x_e(t)$。
\item \textbf{$x_o$ 是奇的}：
$x_o(-t)=\dfrac{x(-t)-x(t)}{2}=-\dfrac{x(t)-x(-t)}{2}=-x_o(t)$。
\item \textbf{加起来是 $x(t)$}：
$x_e(t)+x_o(t)=\dfrac{x(t)+x(-t)+x(t)-x(-t)}{2}=\dfrac{2x(t)}{2}=x(t)$。
\end{enumerate}
三条都成立，所以至少有一种拆法。

\paragraph{先准备一个小引理：既是偶的又是奇的信号，只能是零信号。}
唯一性证明要用到它。这是一个初学时容易意识不到的特性，先单独讲清楚。

\begin{remember}
\textbf{引理。}如果信号 $f(t)$ 既是偶信号又是奇信号，那么 $f$ 是\textbf{零信号}。这里\textbf{零信号}指的是一个信号（函数），它在每一个 $t$ 处的取值都是 $0$，记作 $f\equiv0$。注意区分：「信号是零信号」说的是整条函数，「某个点处取值为 $0$」说的是一个数，引理要证的是前者。
\end{remember}

\textbf{证明。}同一个 $f$ 同时满足两个条件：
\[
\text{偶：}\ f(-t)=f(t),\qquad \text{奇：}\ f(-t)=-f(t).
\]
两个等式的左边是同一个东西 $f(-t)$，所以右边必须相等：
\[
f(t)=-f(t)\ \Longrightarrow\ 2f(t)=0\ \Longrightarrow\ f(t)=0 .
\]
这对\textbf{每一个} $t$ 都成立，所以 $f$ 处处为 $0$。$\blacksquare$

\textbf{用图像理解}：
\begin{itemize}[nosep]
\item 偶信号的图像关于\textbf{纵轴}对称：点 $(t,y)$ 在图上，则 $(-t,y)$ 也在图上；
\item 奇信号的图像关于\textbf{原点}对称：点 $(t,y)$ 在图上，则 $(-t,-y)$ 也在图上。
\end{itemize}
若图像同时满足两种对称，那么点 $(t,y)$ 在图上，就推出 $(-t,y)$ 和 $(-t,-y)$ 都在图上；但一个 $-t$ 处只能有\textbf{一个}函数值，所以 $y=-y$，只能 $y=0$。

\textbf{用一个具体点体会}：设 $f(3)=7$。若 $f$ 是偶的，则 $f(-3)=7$；若 $f$ 又是奇的，则 $f(-3)=-f(3)=-7$。同一个 $f(-3)$ 既是 $7$ 又是 $-7$，矛盾。所以 $f(3)$ 这个\textbf{数}只能是 $0$。对其他任何 $t$ 同理，每个 $t$ 处的取值都是 $0$，合起来就是 $f$ 是零信号。

\textbf{反过来}：零信号确实既是偶的又是奇的（$0=0$，$0=-0$），所以「零信号」是唯一同时具备两种性质的信号。别的信号比如 $\cos t$ 是偶的但不是奇的，$\sin t$ 是奇的但不是偶的，常数 $5$ 是偶的但不是奇的。

\textbf{离散版本}：把 $f(t)$ 换成 $f[n]$，证明逐字相同。

\paragraph{证明二：唯一性（方法 A：两种拆法相减）。}假设有两种拆法
\[
x(t)=e_1(t)+o_1(t)=e_2(t)+o_2(t),
\]
其中 $e_1,e_2$ 是偶的，$o_1,o_2$ 是奇的。移项：
\[
e_1(t)-e_2(t)=o_2(t)-o_1(t).
\]
令 $f(t)$ 等于这个共同的值。
\begin{itemize}[nosep]
\item 作为 $e_1-e_2$，它是偶的：$f(-t)=e_1(-t)-e_2(-t)=e_1(t)-e_2(t)=f(t)$；
\item 作为 $o_2-o_1$，它是奇的：$f(-t)=o_2(-t)-o_1(-t)=-o_2(t)+o_1(t)=-f(t)$。
\end{itemize}
所以 $f$ \textbf{既是偶的又是奇的}，由上面的引理，$f$ 是零信号（$f\equiv0$），即 $e_1-e_2$ 与 $o_2-o_1$ 都是零信号。于是 $e_1=e_2$，$o_1=o_2$：两种拆法其实是同一种。

\paragraph{证明二：唯一性（方法 B：拆法被逼出来了）。}这个方法顺便解释「那个公式是怎么想出来的」。假设已经有某种拆法 $x(t)=e(t)+o(t)$，$e$ 偶、$o$ 奇。把 $t$ 换成 $-t$：
\[
x(-t)=e(-t)+o(-t)=e(t)-o(t).
\]
于是得到两个等式
\[
\begin{cases}x(t)=e(t)+o(t)\\x(-t)=e(t)-o(t)\end{cases}
\ \Longrightarrow\
e(t)=\frac{x(t)+x(-t)}{2},\quad o(t)=\frac{x(t)-x(-t)}{2}
\]
（两式相加得 $e$，相减得 $o$）。也就是说：\textbf{只要拆法存在，$e$ 和 $o$ 就只能是这两个式子，没有别的选择}，唯一性得证。同时它说明课上那个公式不是凭空变出来的，而是从「$x(t)$ 和 $x(-t)$ 两个等式解方程」得到的。

存在性 $+$ 唯一性，定理得证。$\blacksquare$

\paragraph{一个例子。}$x(t)=e^{t}$：$x(-t)=e^{-t}$，所以
\[
x_e(t)=\frac{e^{t}+e^{-t}}{2}=\cosh t,\qquad x_o(t)=\frac{e^{t}-e^{-t}}{2}=\sinh t,
\]
即 $e^t=\cosh t+\sinh t$，偶部是双曲余弦，奇部是双曲正弦。

\emph{（说明：第一节课老师把唯一性留作思考题，说下节课会证；第二节课开头（约 37:11）他说，已经把「奇、偶信号分解是唯一的」的严格证明发在\textbf{课程群}里了，课上没有再讲。所以视频里没有他的证明，以上写法是我（补充）的，其中验算存在性那一段对应他第一节课「自己验证一下」的要求。）}

\textbf{离散的版本}：把 $(t)$ 全换成 $[n]$ 即可：
$x_e[n]=\dfrac{x[n]+x[-n]}{2}$，$x_o[n]=\dfrac{x[n]-x[-n]}{2}$。
老师借此再次强调：连续和离散只是自变量换了写法，道理完全一样，后面很多地方只讲其中之一。

\subsection{能量信号与功率信号}
\textbf{离散}：能量是所有取值的模的平方求和，功率是在有限区间上把能量取平均。
\textbf{连续}：把求和换成积分（老师的说法：「离散的求和与连续的积分是一回事」，因为定积分本来就是切成小段求和）：
\[
E=\lim_{t_0\to-\infty,\ t_1\to+\infty}\int_{t_0}^{t_1}|x(t)|^2\d t,\qquad
P=\lim_{t_0\to-\infty,\ t_1\to+\infty}\frac{1}{t_1-t_0}\int_{t_0}^{t_1}|x(t)|^2\d t
\]
\begin{itemize}[nosep]
\item $E$ 有限 $\Rightarrow$ \textbf{能量信号}（能量有限信号）；
\item $P$ 有限 $\Rightarrow$ \textbf{功率信号}（功率有限信号）。
\end{itemize}
\emph{（补充）功率的表达式我是按板书上的 $t_0,t_1$ 与老师口述的「在有限区间对能量取平均」写的，转写这一段有错字，没有逐项核对。}

%==================================================================
\section{基本的连续时间信号}

\subsection{单位阶跃信号 $u(t)$（书上第 10 页）}
\[
u(t)=\begin{cases}1,&t>0\\0,&t<0\end{cases}
\]
物理含义：电路在 $t=0$ 接入电源，之后电压为 1，之前为 0。

\textbf{$u(0)$ 等于几？} 学生有人答「没定义」，也有答 $0$、$1$、$\tfrac12$。老师说：\textbf{都可以}，取任意值都行，哪怕 $u(0)=-1$ 也可以；不是随口一说，第二章会严格讲，背后有深刻的数学道理。\emph{（补充：根据我对后续内容的理解，这与「一个点的值不影响积分」有关，但课上没展开，我不替老师下结论。）}

\subsection{冲激信号 $\delta(t)$（书上第 12 页）}
\[
\delta(t)=\begin{cases}+\infty,&t=0\\0,&t\neq 0\end{cases}
\qquad\text{且人为规定}\qquad
\int_{-\infty}^{+\infty}\delta(t)\d t=1
\]
直观图像：在 $0$ 处冲上无穷高、再落下来的一个脉冲。书上叫它\textbf{奇异信号}（singular signal，转写里被听成「启艺」）：高中没学过这种函数。严格的数学定义和它在本课里的作用留到第二章。

课件上还写了它想描述的物理现象：时间极短、物理量取值极大，比如雷电、冲击力、电容经小电阻充电。

\subsubsection*{先分清「幅度」和「强度」（课件图上那个箭头旁边的 1 是什么）}
课件把 $\delta(t)$ 画成一根带箭头的竖线，旁边标着 \textbf{1}。这个 1 \textbf{不是高度}：$\delta(t)$ 在 $0$ 处的高度是无穷大，没有一个有限数值。这个 1 是\textbf{强度}（strength，也叫权值），即 $\int_{-\infty}^{+\infty}\delta(t)\d t$，也就是脉冲「面积」的大小。

\begin{remember}
\begin{itemize}[nosep]
\item $\delta(t)$：强度 $=1$（称为\textbf{单位冲激}）。
\item $A\,\delta(t)$：强度 $=A$，即 $\int_{-\infty}^{+\infty}A\delta(t)\d t=A$，$A$ 可以是任何实数（$5$、$-3$ 都行）。
\item 图上画 $A\delta(t)$ 时，箭头旁边标的是 $A$，不是高度。
\end{itemize}
\end{remember}

\textbf{这个「1」从哪来？}它是人为约定的基准，和单位阶跃 $u(t)$（从 $0$ 跳到 \textbf{1}）是同一种做法：先规定一个单位版本，别的大小乘系数就行。可以用一个窄脉冲来看清楚：取宽度 $\Delta$、高度 $h$ 的矩形脉冲，面积 $=\Delta\cdot h$。
\begin{itemize}[nosep]
\item 要面积为 $1$（单位冲激的近似）：$h=\dfrac{1}{\Delta}$；$\Delta=0.1$ 时高 $10$，$\Delta=0.01$ 时高 $100$，$\Delta\to0$ 时高 $\to\infty$，\textbf{但面积始终是 1}。
\item 要面积为 $5$：$h=\dfrac{5}{\Delta}$，极限就是 $5\delta(t)$。
\end{itemize}
所以「幅度趋于无穷、面积却固定」就是 $\delta(t)$ 的全部含义。矩形只是为了好算，三角形、高斯曲线等别的窄脉冲，只要面积是 1、宽度趋于 0，也得到同一个 $\delta(t)$。力学上的类比：冲量 $=$ 力 $\times$ 作用时间，力趋于无穷大、时间趋于零，冲量（面积）是有限值。

\subsubsection*{$\delta(t)=\dfrac{\d u(t)}{\d t}$ 到底怎么理解}
老师给的联系：
\[
u(t)=\int_{-\infty}^{t}\delta(\tau)\d\tau,\qquad \delta(t)=\frac{\d u(t)}{\d t}
\]
\begin{itemize}[nosep]
\item \textbf{积分式好理解}：$t<0$ 时，积分区间 $(-\infty,t)$ 碰不到冲激，结果 $0$；$t>0$ 时，积分区间把整个冲激包了进去，结果是它的强度 $1$。正好就是 $u(t)$。
\item \textbf{求导式要小心}：$u(t)$ 在 $t\neq0$ 处是常数，导数为 $0$；在 $t=0$ 处\textbf{不连续}，按高等数学的规则，不连续点\textbf{普通导数不存在}。所以课件在这里特意写：\emph{由于 $u(t)$ 在跳变点 $t=0$ 处无定义，故上述式不能作为定义式}。也就是说：不能用「阶跃求导」来定义 $\delta(t)$，这个式子只是 $\delta(t)$ 与 $u(t)$ 之间的\textbf{关系}。
\item 工程上放心写这个式子，意思是\textbf{广义导数}（分布意义下的导数）：一个信号在某点有大小为 $A$ 的向上跳变，它的导数在该点带出一个强度为 $A$ 的冲激，即
\[
x(t)=\begin{cases}0,&t<0\\5,&t>0\end{cases}\ \Longrightarrow\ \frac{\d x(t)}{\d t}=5\,\delta(t),\qquad
u(t-2)\ \Longrightarrow\ \frac{\d u(t-2)}{\d t}=\delta(t-2).
\]
对应前面的窄脉冲图像：把阶跃的跳变拉成一小段斜坡，斜坡的斜率是（跳变量）$/\Delta$，斜率就是窄脉冲的高度，面积就是跳变量。
\end{itemize}
老师明确说：现在有个印象就行，学完第二章才会真正理解。

\textbf{关于「不连续所以不能求导」的两个常见误解}（来自我和豆包的对话，我核对后保留）：
\begin{itemize}[nosep]
\item 不是「$\delta$ 不连续所以没法求导」：$\delta(t)$ 根本不是普通函数，「连续、间断」这类逐点的说法都用不上它；它是靠「对别的函数做积分的结果」来定义的（第二章的筛选性质）。
\item 反而是：在广义函数的体系里，间断点也能求导（得到 $\delta$），$\delta$ 本身还能继续求导，$\delta'(t)$ 叫冲激偶，后面会学。
\end{itemize}
\emph{（来源说明：这一小节是我读你和豆包的对话，并对照课件截图后整理的，不是老师课上讲的内容。其中「广义导数」「冲激偶」是课程第二章才系统讲的，这里只作提前说明；我没有查教材，数学上的陈述是我按标准结论核对过的。）}

\subsection{抽样函数 $\Sa(t)$（书上第 14 页，式 1.2.3）}
\[
\Sa(t)=\frac{\sin t}{t}\ (t\neq0),\qquad \Sa(0)=1
\]
\begin{itemize}[nosep]
\item $t=0$ 处不能直接代入（$0/0$），严格定义里 $\Sa(0)=1$，因为 $\lim_{t\to0}\frac{\sin t}{t}=1$。
\item \textbf{偶函数}：$\sin t$ 是奇的，$t$ 是奇的，奇除奇得偶。
\item 零点：$t=\pm\pi,\pm2\pi,\pm3\pi,\dots$。
\item 形状：中间一个主峰（高度 1），两边是振荡衰减的波，因为 $\sin t$ 来回振荡，而分母 $|t|$ 越来越大，把幅度压下去。
\item $\dfrac{\sin\omega t}{t}$（$\omega>0$）的形状（41 分钟老师留给大家先画，43 分钟公布）：形状和 $\Sa(t)$ 一样，但
\begin{itemize}[nosep]
\item 主峰高度：$\displaystyle\lim_{t\to0}\frac{\sin\omega t}{t}=\omega$（老师说用洛必达法则）；
\item 第一个零点：$\omega t=\pi\Rightarrow t=\frac{\pi}{\omega}$，其余零点 $\pm\frac{k\pi}{\omega}$。
\end{itemize}
所以「高度乘 $\omega$，横向往里压缩 $\omega$ 倍」，老师说这是马上要讲的信号变换的伏笔。这个形状同时是后面「面积不变」（积分恒为 $\pi$）的直观来源：高度 $\times\omega$，宽度 $\div\omega$。\emph{（最后一句是我（补充）的联想，不是老师原话。）}
\end{itemize}

\begin{remember}
\[
\int_{-\infty}^{+\infty}\frac{\sin t}{t}\d t=\pi,\qquad
\int_{0}^{+\infty}\frac{\sin t}{t}\d t=\int_{-\infty}^{0}\frac{\sin t}{t}\d t=\frac{\pi}{2},\qquad
\int_{-\infty}^{+\infty}\frac{\sin\omega t}{t}\d t=\pi\ (\omega>0)
\]
\end{remember}
老师说：这三个公式第三章推傅里叶变换要用，必须背。

\subsubsection*{三个公式之间的关系}
\begin{enumerate}
\item \textbf{整条线 $\leftrightarrow$ 半条线}：$\frac{\sin t}{t}$ 是偶函数，图像关于纵轴对称，所以 $0$ 右边和 $0$ 左边的积分相等，各占一半：$\pi/2$。
\item \textbf{$\omega$ 的版本}：令 $s=\omega t$（$\omega>0$ 保证积分上下限不变号），则 $t=s/\omega$，$\d t=\d s/\omega$：
\[
\int_{-\infty}^{+\infty}\frac{\sin\omega t}{t}\d t
=\int_{-\infty}^{+\infty}\frac{\sin s}{s/\omega}\cdot\frac{\d s}{\omega}
=\int_{-\infty}^{+\infty}\frac{\sin s}{s}\d s=\pi
\]
$\omega$ 约掉了。（这是老师 33:30--35:00 在黑板上做的变量代换；40:44 他又说「这个变种上次课证明过」。）
\end{enumerate}

\subsubsection*{证明 $\displaystyle\int_0^{+\infty}\frac{\sin t}{t}\d t=\frac{\pi}{2}$（板书约 43--54 分钟，参数求导法）}
\begin{note}
这一段是 44:16--53:00，老师说这是「超越书上内容」的补充：书上的证明要用复变函数，他给的是不用复变函数的初等证明；目的有两个，一是严谨，二是让大家记住那三个公式（它们在第三章开始不断用到）。下面的步骤与理由按老师的讲解整理，我对照了板书；各步的措辞是我的，数学上我自己代入验算过（见第 3 步）。
\end{note}

\textbf{思路}：直接积 $\frac{\sin t}{t}$ 积不出来（原函数不是初等函数）。于是在积分里人为添一个衰减因子 $e^{-at}$，把问题变成「含参数 $a$ 的积分」，对 $a$ 求导，导数是好积的，最后让 $a\to0$ 把因子去掉。

\textbf{第 1 步：设}
\[
I(a)=\int_0^{+\infty}\frac{\sin t}{t}\,e^{-at}\d t\qquad(a\ge 0)
\]
想要的是 $I(0)$。

\textbf{第 2 步：对 $a$ 求导}（把求导放进积分号，$\frac{\partial}{\partial a}e^{-at}=-t\,e^{-at}$，正好把分母的 $t$ 消掉）：
\[
\frac{\d I}{\d a}=-\int_0^{+\infty}\sin t\;e^{-at}\d t
\]

\textbf{第 3 步：用欧拉公式}(Euler's formula) $\sin t=\dfrac{e^{jt}-e^{-jt}}{2j}$：
\begin{align*}
\frac{\d I}{\d a}
&=-\frac{1}{2j}\int_0^{+\infty}\bigl(e^{jt}-e^{-jt}\bigr)e^{-at}\d t\\
&=-\frac{1}{2j}\Bigl[\int_0^{+\infty}e^{-(a-j)t}\d t-\int_0^{+\infty}e^{-(a+j)t}\d t\Bigr]\\
&=-\frac{1}{2j}\Bigl[\frac{1}{a-j}-\frac{1}{a+j}\Bigr]
=-\frac{1}{2j}\cdot\frac{2j}{a^2+1}
=-\frac{1}{a^2+1}
\end{align*}
（这里 $a>0$ 保证 $e^{-(a\mp j)t}\to0$，积分才收敛。老师口头把负号并进括号里，写成 $\frac{1}{2j}\bigl[\frac{1}{a+j}-\frac{1}{a-j}\bigr]$，通分后分母 $a^2-j^2=a^2+1$、分子 $-2j$，结果一样；他说「积分学得好的一眼能看出」，并让大家课后自己验证这一步。）

\textbf{第 4 步：积分求 $I(a)$}：
\[
I(a)=-\arctan a+C
\]

\textbf{第 5 步：用 $a\to+\infty$ 定常数 $C$}。$a$ 很大时 $e^{-at}$ 把整个被积函数压成 $0$，所以 $I(+\infty)=0$：
\[
0=I(+\infty)=-\arctan(+\infty)+C=-\frac{\pi}{2}+C\ \Longrightarrow\ C=\frac{\pi}{2}
\]

\textbf{第 6 步：令 $a\to 0$}：
\[
I(0)=\int_0^{+\infty}\frac{\sin t}{t}\d t=-\arctan 0+\frac{\pi}{2}=\frac{\pi}{2}\qquad\blacksquare
\]

再由偶函数得 $\int_{-\infty}^{+\infty}\frac{\sin t}{t}\d t=2\cdot\frac\pi2=\pi$。

\emph{（补充，严格性：第 2 步「求导进积分号」和第 6 步「$a\to0^{+}$ 时 $I(a)\to I(0)$」严格来说需要论证收敛性，重转写里老师没有展开，他是直接把 $a=0$ 代入的，这一点是我看出来的，不是他说的。）}\par\smallskip\textbf{老师在第 5 步的原话要点}：$a\to+\infty$ 时 $e^{-at}$ 趋于 0，被积函数每一点都趋于 0，所以 $I(+\infty)=0$；第 6 步「最后一步很有趣」：令 $a=0$，$e^{-0\cdot t}=1$，$I(0)$ 就是要求的积分。

%==================================================================
\section{基本的离散时间信号（书上第 15 页起）}

\subsection{单位脉冲序列 $\delta[n]$}
\[
\delta[n]=\begin{cases}1,&n=0\\0,&n\neq0\end{cases}
\]
书上叫它单位脉冲序列(unit impulse sequence)。

\subsection{单位阶跃序列 $u[n]$}
\[
u[n]=\begin{cases}1,&n\ge0\\0,&n<0\end{cases}
\]

\subsection{连续与离散的四个信号对照}
\begin{center}
\begin{tabular}{lcc}
\toprule
 & 连续（圆括号） & 离散（方括号）\\
\midrule
冲激 & $\delta(t)$：奇异，$t=0$ 处无穷大 & $\delta[n]$：普通，$n=0$ 处等于 1\\
阶跃 & $u(t)$：$t>0$ 为 1，$u(0)$ 任意 & $u[n]$：$n\ge0$ 为 1\\
\bottomrule
\end{tabular}
\end{center}
老师特别强调：\textbf{$\delta(t)$ 和 $\delta[n]$ 不是一回事}，常被搞混。$\delta[n]$ 是正常得不得了的信号（零处为 1，别处为 0）；$\delta(t)$ 是高中没学过的奇怪信号，第二章才会严格讲。区分的办法就是前面「圆括号连续、方括号离散」那条规则。

%==================================================================
\section{本节要记住的清单}
\begin{enumerate}[nosep]
\item $x(t)$ 连续，$x[n]$ 离散；没有 $x[t]$ 与 $x(n)$。
\item 偶 $+$ 奇分解：$x_e=\frac{x(t)+x(-t)}2$，$x_o=\frac{x(t)-x(-t)}2$，唯一。
\item 能量信号：$E<\infty$；功率信号：$P<\infty$。
\item $u(t)$：$t>0$ 为 1、$t<0$ 为 0，$u(0)$ 可任取；$\delta(t)=\frac{\d u}{\d t}$，$\int\delta=1$。
\item $\Sa(t)=\frac{\sin t}{t}$：偶函数，$\Sa(0)=1$，零点 $\pm k\pi$。
\item $\int_{-\infty}^{\infty}\Sa(t)\d t=\pi$；$\int_0^{\infty}\Sa(t)\d t=\frac\pi2$；$\int_{-\infty}^{\infty}\frac{\sin\omega t}{t}\d t=\pi\ (\omega>0)$。
\item $\delta[n]$、$u[n]$ 是离散里的对应物；$\delta[n]\neq\delta(t)$。
\end{enumerate}

\end{document}
